% Edition 19: first essay, 8.14--8.18.
% French source: physical pages 46 (Combinant 8.13) through 50
% (N.B. immediately before Il reste a expliciter).
% English editorial source: physical pages 32--34, before post-8.18 discussion.
% Replaces the old DOT_IMPORT 08, 09, 10 only. Old sources remain immutable.
\begingroup
\def\PD{P^\vee}
\def\depth{\operatorname{depth}}
\def\codim{\operatorname{codim}}
\source{法文手稿物理第46--50页；英编稿物理第32--34页。第十九版回校}

结合8.13、8.5和8.6，按通常方式可得如下推论。\transnote{这是作者的承接句。法文第46页以手写补入“及8.6”。以下8.14保留法稿修订层的充要措辞；8.13的充分性本身并不能代替各必要方向的证明。本次源页辨读、必要的前提核对和未完整重证的范围见核查报告“D19-8.14--8.18”。}

\egas{8.14}{推论（法稿修订层）}
假设 $f\colon X\to P$ 非分歧，且 $X$ 没有孤立点。\transnote{法文第46页的这两个假设可读；其下一行左侧另有一个孤立的打字字母 $n$，没有形成数学条件。本版将该残字登记在校记中，不把它接成一条缺失的附加假设。英编稿在这里留下第49注，原文为“询问 A.G.”，仍予保留。}\textsuperscript{,}\ednote{询问 A.G.。〔英编稿脚注49。〕}
\begin{enumerate}
\item[(a)] 假设 $X$ 在 $k$ 上可分。为使那些满足 $\varphi_\xi$ 为 $X_{k(\xi)}$-正则、且 $Y_\xi$ 可分的点 $\xi\in\PD$ 所成集合，其补集的余维至少为 $2$，充要条件是：$X$ 的每个一维不可约分支 $X_i$ 都由 $X$ 的光滑点组成，并且相对于 $f$ 是“直纹的”；此外，对于 $X$ 的每个满足 $\dim_xX\ge2$ 的闭点 $x$，都有\transnote{“可分”沿用本篇旧用语，指几何约化，见本篇3.3以及第十八版第91批相应说明，不是分离性或光滑性的同义词。“直纹的”沿用8.13中切关联像的余维工作含义，不是通常的直纹簇。法稿第46页的旧条件被框划，现条件从第47页首接入；第47页只写一次 $\operatorname{prof}_xX\ge2$。英编稿重复的深度因子与误排的 $\phi_{xi}$ 均据此订正。}
\[
 \depth\mathcal O_{X,x}\ge2.
\]
（若 $X$ 几何正规，且所有不可约分支的维数均至少为 $2$，这些条件便自动满足。）
\item[(b)] 假设 $X$ 几何正规。为使那些满足 $\varphi_\xi$ 为 $X_{k(\xi)}$-正则、且 $Y_\xi$ 几何正规的点 $\xi\in\PD$ 所成集合，其补集的余维至少为 $2$，充要条件是：$X$ 的每个维数不超过 $2$ 的不可约分支 $X_i$ 都由 $X$ 的光滑点组成，并且相对于 $f$ 是“直纹的”；此外，对于 $X$ 的每个满足 $\dim_xX\ge3$ 的闭点 $x$，都有\transnote{法文第47页(b)的指标为“不超过2”，末式为深度至少3。底本在本条之后还有约三行经密集覆打删除的文字；未能恢复的字词不补造成另一条假设，校记另存。涉及几何性质时，仍须连同前文的几何换底约定阅读。}
\[
 \depth\mathcal O_{X,x}\ge3.
\]
\end{enumerate}

\egas{8.15}{备注}
在8.6、8.13和8.14中，我们对 $X$ 作出 $(S_k)$（分别地，$(R_k)$；又分别地，可分、几何正规）这一假设，而希望将同一性质作为超平面截面的结论，至多排除一个余维至少为 $2$ 的例外参数集。

这并不限制一般性；说实话，省去这一预先假设也不会多费事。因为借助第3、4段和5.12的结果，可以立即看出：若 $X$ 不满足所讨论的假设，那么，由第5段，存在一个使该假设失效的闭点 $x$；对每个满足以下条件的 $\xi$：\transnote{这里按法文第47页的 \emph{par. 3.4.} 保留“第3、4段”的读法，5.12为手写补引。该句为作者草稿的论证提议；前述条目的有限型、正则截面等前提仍须保留，不能由这句话无条件删去。}
\begin{enumerate}
\item $\varphi_\xi$ 为 $X_{k(\xi)}$-正则（这个条件只排除一个余维为〔右缘缺损，数值未辨〕的集合）；
\item $x\in Y_\xi$（这个条件给出一个余维恰为 $1$ 的集合），
\end{enumerate}
$Y_\xi$ 在 $x$ 处也都不满足上述假设。因此，例外集合 $Z\subset\PD$ 的余维是 $1$，而不是 $2$。

（在 $(R_k)$ 的情形，我也许说得有些过头了；在那里还需要某些条件：$(S_1)$，也许还需要等维性……。）\transnote{法文第47页明确为 $(R_k)$，不是英编稿的 $E_k$；这项作者保留意见是原文的组成部分。括号前非正则截面例外集的余维数值处在扫描右缘外，仍不能恢复。}

\subsection*{作者补记：重新考虑（Remords）}
作者建议把8.13、8.14的非分歧假设改为只在 $X$ 的光滑点处要求；他还建议，对充分性方向，只要求在补集余维至少为 $k+1$ 的某个开集 $X'$ 上非分歧。\transnote{此为法文第48页原补记的完整两项建议。第二项不能据所写余维界直接作为一般定理采用；下面8.16的原弱假设恰受此影响，核查报告给出光滑射影三维概形上的反例。这里保留作者提议的身份，不据此删除8.13、8.14正文中的全局非分歧假设。第一项也未在本次重证为任意情形下的替换规则。}

\egas{8.16}{命题（译校采用的加强假设版本）}
假设 $f\colon X\to P$ 非分歧，$X$ 几何正规，且在每个闭点处的深度至少为 $3$；最后，假设 $X$ 几何整且在 $k$ 上固有。那么，使得 $Y_\xi$ 几何正规、几何整且维数等于 $\dim X-1$ 的点 $\xi\in\PD$ 所成集合（是可构造的，并且）其补集的余维至少为 $2$。\transnote{法文第48页及英编稿第33页原只要求 $f$ 在某开集上非分歧，且该开集的补集余维至少为2。这个原条件不足：核查报告“D19-8.16：余维二收缩反例”给出满足原文全部假设、而坏参数包含一个超平面的例子。正文仅把这一项加强为全局非分歧，完整保留原有深度、几何正规、几何整和固有假设；不是声称恢复了作者原有的正确定稿。原陈述及全部原证明另存报告。}

\begin{proof}[加强假设下的证明]
闭点深度条件蕴涵 $\dim X\ge3$。扩域到 $\overline{k}$ 时，几何闭点映到原概形的闭点；原闭点局部环中长度为3的正则列，在平坦局部扩张到相应几何闭点局部环后仍为正则列，因此所需的几何闭点深度条件仍成立。由8.6，除了余维至少为2的参数集以外，正则截面满足 $(S_2)$；由8.13(b)取下标 $n=1$，同样只需排除余维至少为2的参数集，便使截面几何满足 $(R_1)$。这里没有维数不超过2的不可约分支，而几何整性也排除了零维伴随点。因此这些截面几何正规且维数为 $\dim X-1$。这说明作者在原证明中对8.14(b)的应用，在本版加强的前提下成立。

又因 $f$ 非分歧，$\dim f(X)=\dim X\ge3$。对几何整的固有概形应用6.1，可知所有这些 $Y_\xi$ 都几何连通。几何正规且非空的有限型概形，几何整当且仅当几何连通，于是得到结论。\transnote{前一段将原证明省略的 $(S_2)$、$(R_1)$ 和伴随点条件明确展开，属于译校补证。连通性这一步保留原作者的论证；其对像维数和固有性的使用在此明列。原弱假设版本不能仅靠重复该证明获得验证。}
\end{proof}

\egas{8.17}{备注及作者未定稿的补充}
\begin{enumerate}
\item[(a)] 8.16的假设蕴涵 $\dim X\ge3$。单有“$X$ 几何不可约且 $\dim f(X)\ge3$”，不足以使几何不可约超平面截面的坏参数集具有至少为2的余维；坏参数集甚至可以包含一个超平面。\transnote{原稿猜测，去掉正规性和非分歧性后，几何不可约及 $\dim f(X)\ge3$ 也许仍足够，并接着断言“无论如何，可以证明这个补集不包含任何超平面（见下文）”。这两项宽泛说法均被 $f\colon\operatorname{Bl}_p(\mathbb P^3)\to\mathbb P^3$ 反驳：每个经过 $p$ 的超平面的总变换同时含例外除子和严格变换，故可约；这些参数组成对偶空间中的一个超平面。正文写明正确的否定结论，原语气和原断言完整保存在本注及报告。此反例不满足8.18的有限集合外非分歧假设。}
\item[(b)] 若去掉闭点 $x$ 处 $\depth\mathcal O_{X,x}\ge3$ 的假设，8.16的结论便不成立。例如，对 $\mathbb P^3$ 中的一个非奇异二次曲面 $X$，结论就不成立。与该二次曲面相切的超平面截面是可约的，事实上由两条相交直线组成；这些切超平面形成一个二维族，因而在 $\PD$ 中的余维为1（事实上，该二次曲面的对偶，是对偶空间中由对偶形式给出的二次曲面……）。\transnote{法文第48页清楚写有 \emph{hyperplans tangents}，英编稿在此留出的缺文因而可以恢复。原稿在句法上把“可约”直接接在切超平面之后，本版明说为相应的截面。法稿参数空间已有对偶记号，英稿此处的 $P$ 据此校正。关于“对偶形式”的括注按作者原话保留，不据此扩张成任意特征下的无条件线性代数断言；报告用特征零的光滑分裂二次曲面核实本例。}\textsuperscript{,}\ednote{询问 A.G.。〔英编稿脚注50；所附“相切”文字现已由法稿辨明。〕}
\end{enumerate}

对于射影空间中的非奇异曲面，这种情形却应当视为例外情形，参看下一编号小节。

\paragraph{(c)（作者手写补充）}
假设 $X$ 是光滑的二维概形，几何整且在 $k$ 上固有，并且 $f$ 是浸入。由6.1、8.8及8.14可得：如果 $g\colon Y^{\mathrm{sing}}\to\PD$ “不是泛有限且不可分的”，那么，那些满足 $Y_\xi$ 在 $k(\xi)$ 上可分、并且至多有两个几何不可约分支的 $\xi\in\PD$ 所成集合，其补集的余维至少为2。\transnote{法文第49页顶端手写段标作(c)；上述维数、光滑、几何整、固有和浸入假设及“至多两个几何不可约分支”均可辨，8.14是8.8旁的补引。原条件逐字为 \emph{n'est pas gén. fini et inséparable}；“不”的逻辑作用范围没有另行写明。本版保留这一歧义及原稿未定稿的身份，并且8.8缺失起首的问题仍在，不能将此段当作新确立的普遍充分条件。英编稿的若干缺文在这里已得到恢复，不再一概标作不可辨。}

现在我们要更仔细地考察曲面的情形（从超平面截面的不可约性这一观点看，曲线的情形显然不成为问题）。

（注意：我惊恐地发现，按我一直使用“直纹的”一词时所取的意义，二次曲面竟没有资格称为直纹的！这似乎与前人的用法不一致，因而有必要为这里使用的概念另造一个更恰当的词。）\transnote{这是作者对临时术语的提醒，不是否定通常意义下光滑分裂二次曲面的直纹结构。}

\egas{8.18}{命题（译校采用的充分条件版本）}
设 $k$ 为代数闭域，$X$ 为在 $k$ 上固有的正规整概形，$d=\dim X\ge2$。设 $T\subset X$ 为有限闭子集，使 $X'=X\setminus T$ 光滑，且 $f|_{X'}$ 非分歧。记
\[
 D=\overline{g(Y'^{\mathrm{sing}})}\subset\PD.
\]
假设下列两个条件成立：\transnote{原稿包括“不要求 $X$ 正规”的分别情形，并称原(a)、(b)为充要条件；它的首句仍印“维数1”，而条件及证明已手改为 $d-1$。更实质的问题是原(a)、(b)只要求截面不可约，这允许二倍直线；二次 Veronese 曲面的例子已反驳原结论，详见报告“D19-8.18：原条件的反例与充分版本”。本版明确采用正规情形和下面加强的几何整条件，结论只称充分，不把它冒作原一般命题的等价恢复。原完整陈述、全部作者证明、被划去证明及末尾附记均在报告中逐段保存。}
\begin{enumerate}
\item[(a)] 对每个 $x\in T$，存在一个经过 $x$、几何整且维数为 $d-1$ 的超平面截面 $Y_\xi$，其中 $\xi\in\PD(k)$。
\item[(b)] 或者 $D$ 的余维至少为2（包括 $D=\varnothing$ 的情形），或者存在 $X'$ 的一个超平面截面 $Y'_{\xi_0}$，其中 $\xi_0\in\PD(k)$，它奇异、几何整且维数为 $d-1$。\transnote{法文第49页(b)的末行明确为 \emph{singulière et irréductible}：“奇异且不可约”。它不是英编稿疑作的“(non???) singular”。正文的“几何整”是译校加强，而“奇异”是法稿的实际读法。原文的第一种情况用临时术语“直纹的”；这里以切关联像闭包的余维明确本版本采用的工作含义。}\textsuperscript{,}\ednote{字迹不清，询问 A.G.。〔英编稿脚注51；其中奇异／非奇异的疑读已由法稿辨明。〕}
\end{enumerate}
那么，使得 $Y_\xi$ 几何整且维数为 $d-1$ 的点 $\xi\in\PD$ 所成集合，其补集的余维至少为2。

\begin{proof}[译校补证（非原文）]
令 $U$ 为满足结论中两项性质的参数集。先限制到截面方程正则的开集；通用截面族在其上固有、平坦且有限表示，故几何整纤维的参数构成开集。因此 $U$ 在 $\PD$ 中开。

对每个 $x\in T$，令 $H_x$ 为经过 $f(x)$ 的超平面所成的参数超平面，并置
\[
 Z=D\cup\bigcup_{x\in T}H_x.
\]
在 $\PD\setminus Z$ 上，截面避开 $T$，且在 $X'$ 上光滑、维数为 $d-1$。由于 $f|_{X'}$ 非分歧，$\dim f(X)=d\ge2$；由6.1，截面几何连通，因而几何整。所以 $\PD\setminus Z\subset U$。

切关联概形 $Y'^{\mathrm{sing}}$ 非空时是不可约的：它是光滑不可约概形 $X'$ 上的射影丛。因此 $D$ 或为空，或为不可约真闭集。$Z$ 的余维一分支仅可能是 $H_x$ 以及余维为1时的 $D$。条件(a)使每个 $H_x$ 与 $U$ 相交。

若 $D$ 的余维为1，条件(b)给出的 $Y'_{\xi_0}$ 几何整。原完整截面 $Y_{\xi_0}$ 是正规概形 $X$ 上的有效 Cartier 除子，故满足 $(S_1)$；它的所有不可约分支维数均为 $d-1\ge1$，从而没有分支包含于有限集 $T$。于是 $Y'_{\xi_0}$ 的几何整性蕴涵 $Y_{\xi_0}$ 的几何整性，并且完整截面确为该开截面的闭包。又因该开截面奇异，$\xi_0\in D\cap U$。

故 $Z$ 的每个余维一分支均与开集 $U$ 相交。闭集 $\PD\setminus U$ 包含于 $Z$，且不含 $Z$ 的任何余维一分支，所以其余维至少为2。\transnote{本补证使用几何整纤维的开性，而不把仅有“不可约”的纤维误作开条件；也不以一个非约化但不可约的特殊截面代替 $D\cap U$ 中的点。这里已经假设 $X$ 正规，故不需要采用原稿略写且未核全的正规化归约。相关局部环论证及原稿在像闭包上的正确横线位置见报告。}
\end{proof}

\paragraph{作者附记的保存位置.}
法文第50页末的附记确实说：在“分别地”的情形，8.18包含在8.16中，除了 $d=2$；对不带“分别地”的情形，也许应当不要求 $d=2$，而只要求 $d\ge2$。\transnote{末尾“只要求”可从法文手迹恢复，英稿的 \emph{and not only} 不能据以保留反向读法。该附记是作者对原草稿的编写意见；它既没有补入8.16的闭点深度条件，也没有证明正规化归约，不能把它当作本版两个加强版本之间的包含关系。完整原附记与原证明连续存于核查报告。}
\par
\endgroup
